
\subsection{A useful text signal still needs an intervention argument}
\label{rel66:text-signals}

The distinction between a label and an action also governs claims that a
language model has discovered causes inside case notes. A bank may find that
notes mentioning a payment reminder have higher predicted repayment. Such a
signal can help an analyst search records or formulate a hypothesis. To decide
whether to send a reminder, however, the bank needs the repayment contrast
under an implementable reminder policy and its alternative, using information
available before that choice. A word recorded after the action cannot silently
become a baseline treatment assignment or an admissible adjustment variable.

Three pharmacovigilance studies provide a concrete comparison.
\citet{rel66inferbert} trains a classifier on sentences constructed from
commercially curated FAERS fields. Its two reported tasks concern acute liver
failure and tramadol-related death. Equations 2--6 form groups according to
whether an observed field contains a term, then compare the classifier outputs
in those groups. The paper calls this operation do-calculus. In the inspected
\texttt{causal\_tree} and \texttt{compute\_gap} routines, the predictions are
read from a saved file and partitioned by feature membership; the procedure
does not replace a structural equation or evaluate both actions for the same
unit. Recursive restriction to records containing an earlier term does not
itself supply an adjustment argument. \citep{rel66inferbertcode}

To state what this computes, let $W_i$ denote the recorded text or fields,
$H_i=1$ indicate a selected term's presence, and $f(W_i)$ be a fixed predicted
endpoint probability. On the retained records, the root comparison is
\begin{equation}
 \widehat\Delta_H=
 \frac{\sum_i H_i f(W_i)}{\sum_i H_i}
 -\frac{\sum_i(1-H_i)f(W_i)}{\sum_i(1-H_i)}.
 \label{eq:rel66-score-gap}
\end{equation}
This expression assumes nonempty groups; the code further restricts to
nonmissing fields and groups exceeding its support threshold. Its normal-score
denominator uses the sum of the two within-group sample variances divided by
their respective counts. That calculation describes uncertainty about a
particular score contrast only under its own sampling and fitting assumptions.
It does not establish exchangeability, resolve reporting selection or account
automatically for learning and repeatedly selecting the highest-scoring terms.
The published paper describes adjusted significance levels, while the selected
routine visibly applies an unadjusted one-sided threshold of 0.05. The local
source inspection therefore cannot certify the paper's complete multiplicity
procedure or reconstruct its reported results. \citep{rel66inferbert,rel66inferbertcode}

\citet{rel66deepcausality} extends this design to free text from LiverTox,
using named-entity recognition to propose terms. Its label is inherited from a
drug-level likelihood category: the extracted sentences receive that drug's
positive or negative label. The reported split places 14,361 sentences into
training and test sets in a 9:1 ratio; the account selects its checkpoint using
the test-set classification results. This does not supply an untouched
drug-level evaluation. It leaves both selection optimism and within-drug
dependence to be assessed; no numerical penalty is inferred here. The reported
18 clinically consistent terms out of 20 retained terms measure agreement with
domain knowledge. They do not measure intervention-effect bias or interval
coverage. The paper itself states that the method cannot classify the selected
factors into confounders and colliders. In the pinned notebook, the same score
partition is visible and the significance-filter lines are commented out.
These are reasons to specify the exact implementation and evaluation,
including useful hypothesis-ranking behavior, before adopting its causal
interpretation. \citep{rel66deepcausalitycode}

The subsequent comparison by \citet{rel66pharma} retains both the structured
and free-text tasks. Its domain-expert assessment asks whether enriched terms
are relevant to liver injury. The stated positive predictive quantity is
$\mathrm{TP}/(\mathrm{TP}+\mathrm{FP})$; the narrative reports 0.95 for the
task-adapted DILI-BERT and 0.20 for RoBERTa. Those numbers concern this selected
term-review task. They cannot be read as probabilities that a banking
intervention estimate is correct. The method paragraph cites references 9--10,
whereas the bibliography identifies InferBERT and DeepCausality as 17--18;
the comparison here follows the actual primary methods, rather than treating
the mismatched reference numbers as a new derivation. The source's model-size
and domain-adaptation comparisons remain useful experimental leads within
their reported task and cohort.

An exact authored example isolates the missing step without imitating any of
these experiments. Let unrecorded customer readiness $U$ be Bernoulli with
probability $1/2$, and let a historical action $A$ have probabilities
$\Pr(A=1\mid U=1)=4/5$ and $\Pr(A=1\mid U=0)=1/5$. Suppose repayment $Y$ has
conditional probability $1/10+3U/5$ and is independent of $A$ given $U$.
Independent uniform random variables can implement these two Bernoulli draws,
so intervening on $A$ leaves the repayment equation unchanged. A note records
$A$ exactly. Bayes' rule gives
\begin{equation}
 E[Y\mid A=1]=\frac{29}{50},\qquad
 E[Y\mid A=0]=\frac{11}{50},\qquad
 E[Y(1)-Y(0)]=0.
 \label{eq:rel66-confounding}
\end{equation}
Even the exactly calibrated predictor $f(A)=E[Y\mid A]$ therefore produces
a term-group gap of $9/25=0.36$. For a stipulated USD~100 value per repayment
and USD~10 action cost per account, treating the gap as an effect implies
USD~26 net value; the stipulated intervention law gives minus USD~10. This
is a finite probability construction, not a bank estimate or a replication of
the source studies. Better prediction of this same observational law cannot
choose the causal interpretation.

\subsection{A conditional route from text to adjustment}

Text can nevertheless contain useful adjustment information. The route in
\citet{rel66veitch} begins by assuming that the original text suffices for
the specified causal adjustment. It then asks whether a lower-dimensional
embedding retains the outcome or propensity information needed for estimation.
This addresses loss of already sufficient information during compression; it
does not establish that an arbitrary case note captured all common causes.
Its semi-synthetic experiments test the proposed representations under stated
generators. Its consistency result also requires sufficiently accurate nuisance
models and a well-behaved estimator, with bounded outcomes and overlap given
as sufficient conditions in the discussion. Predictive language performance
alone does not verify these requirements in a bank.

For a single-stage banking use, take $W$ to be text available before a defined
action $A$, let $Z=\lambda(W)$ be its representation, and let $Y$ be the later
outcome. State consistency, conditional exchangeability given $W$, and overlap.
If $Z$ preserves the required adjustment information, define
$m_a(z)=E[Y\mid A=a,Z=z]$ and $e(z)=\Pr(A=1\mid Z=z)$. The target population
then determines the averaging law:
\begin{align}
 \tau_{\mathrm{ATT}}
 &=E[m_1(Z)-m_0(Z)\mid A=1]
   =\frac{E[e(Z)\{m_1(Z)-m_0(Z)\}]}{\Pr(A=1)},\\
 \tau_{\mathrm{ATE}}
 &=E[m_1(Z)-m_0(Z)].
 \label{eq:rel66-target-weights}
\end{align}
The equality for ATT follows from Bayes' rule after identification. It shows
why a representation model and a correct conditional outcome contrast are
insufficient if the final estimator averages over the wrong population.
In the inspected proceedings version, equation 2.1 conditions the outer
expectation on treatment, and equations 2.2--2.3 retain propensity or treated
weights. The later displayed equation 3.1 and the final display in Theorem 3.1
omit those weights or that conditioning while discussing the same target.
We retain the explicit ATT definition and distinguish the displayed formulas;
we do not treat this notation mismatch as a theorem that ATT equals ATE.
The separate appendix proof is not contained in the inspected ten-page
proceedings copy. \citep{rel66veitch}

The distinction can reverse a decision even when all confounding is observed.
Let $Z$ have two equally prevalent strata, with propensities $1/5$ and $4/5$,
and conditional repayment effects zero and $3/10$. One feasible law sets
$m_0(z)=1/5$ in both strata and $m_1(z)=1/5$ and $1/2$, respectively.
Then $\Pr(A=1)=1/2$, the treated population places weights $1/5$ and $4/5$
on the two strata, and
\begin{equation}
 \tau_{\mathrm{ATE}}=\frac{3}{20}=0.15,\qquad
 \tau_{\mathrm{ATT}}=\frac{6}{25}=0.24.
\end{equation}
At USD~100 per repayment and USD~20 per action, the respective net contrasts
are minus USD~5 and plus USD~4 per account in their respective populations.
Neither is a universal policy value: expanding eligibility beyond historical
recipients changes the population and may change support. The example holds
the conditional effects fixed so that the consequence of weighting is visible.

The proposed service should consequently distinguish three outputs: a
term-ranking diagnostic, a measured pretreatment variable, and an identified
effect under a declared adjustment argument. A credible text-assisted study
retains source spans and availability times, evaluates extraction on held-out
customers and periods, and propagates measurement and nuisance-estimation
uncertainty using a sampling design appropriate to repeated notes. It compares
the text representation with structured covariates and simpler text summaries,
reports overlap and sensitivity to omitted information, and uses an untouched
evaluation set after model selection. If the adjustment argument fails,
hypothesis ranking or a randomized prospective action comparison remains a
substantive alternative. These requirements preserve the useful text-processing
contribution without turning its classifier score into a causal assurance.
